% =====================================================================
%  RESUMEN Y ABSTRACT
%  La plantilla fija 150-250 palabras y pide objetivo, metodo, resultados y
%  discusion. Las palabras clave son obligatorias en los dos idiomas.
%
%  REGISTRO: un resumen comprime, asi que sus frases salen mas largas que las
%  del cuerpo. Aun asi se mantienen por debajo de las 30 palabras, que es el
%  limite que el resto del documento respeta.
% =====================================================================
\ustaseccion{Resumen}

Orientarse en un edificio de varios pisos es difícil para quien no lo conoce, y
el cambio de nivel corta la continuidad del recorrido. Las soluciones robóticas
existentes resuelven ese salto de dos maneras: dotando al robot de mecanismos
para subir escaleras, o conectándolo al ascensor. Las dos son caras y obligan a
intervenir el edificio.

Este trabajo propone y evalúa una alternativa. Dos agentes, robots móviles terrestres, uno por planta, se reparten el guiado de una persona mediante un protocolo de relevo
en la zona de escaleras. Ningún agente cruza de nivel. Un coordinador central
asigna el agente según el nivel del origen y del destino. El usuario pide el
servicio desde una interfaz web que funciona en cualquier teléfono, sin instalar
nada y sin salida a internet.

La evaluación siguió un protocolo escrito antes de instrumentar el sistema, con
cuatro métricas de definición operativa. En simulación, treinta misiones
sorteadas dieron una campaña válida sin descartes, que permite afirmar con 95~\%
de confianza que la tasa de éxito supera el 70~\%.

De diez misiones sorteadas con los agentes físicos, uno por piso, nueve se
completaron con relevo, con un tiempo total mediano de 87,6~s y la
llegada a menos de 0,5~m del destino según la localización de cada agente. La
décima se interrumpió porque el computador del coordinador se quedó sin batería.
La tracción no responde por debajo de 0,40~m/s, y eso limita la precisión de llegada.

\vspace{1em}
\textbf{Palabras clave:} robótica móvil colaborativa; sistemas multi-robot;
navegación en interiores; protocolo de relevo; interacción humano-robot; ROS~2.

\ustaseccion{Abstract}

Finding one's way inside a multi-storey building is hard for first-time
visitors, and changing floors breaks the continuity of the route. Existing
robotic solutions address that gap in two ways: by equipping the robot to climb
stairs, or by interfacing it with the building lift. Both are expensive and
require modifying the infrastructure.

This work proposes and evaluates an alternative. Two agents, ground mobile robots, one per floor, share the task of guiding a person through a handover protocol at the
stairwell. Neither agent crosses between levels. A central coordinator assigns
the agent according to the floor of the origin and the destination. The user
requests the service from a web interface that runs on any phone, with no
installation and no internet access.

The evaluation followed a protocol written before instrumenting the system, with
four operationally defined metrics. In simulation, thirty randomised missions
gave a valid campaign with no discards, which supports, with 95~\% confidence,
the claim that the success rate exceeds 70~\%.

Of ten randomised missions with the physical agents, one per floor, nine were
completed with a handover, with a median total time
of 87.6~s and arrival within 0.5~m of the goal according to each agent's own localisation. The tenth was interrupted when the coordinator's computer ran out
of battery. Traction does not respond below 0.40~m/s, which limits arrival accuracy.

\vspace{1em}
\textbf{Keywords:} collaborative mobile robotics; multi-robot systems; indoor
navigation; handover protocol; human-robot interaction; ROS~2.
